\section{Performance Within Operating Regions}\label{sec:within}
The operating-region names identify FIFO-referenced arrival-rate settings,
not independently measured physical saturation states. Each source has
40 windows and every primary policy is observed on each window in each
region. Comparisons within a source/region use matched windows; a policy
is not calibrated separately to obtain a favorable operating point.

The central quantity is relative order under ANWG, the fraction of arrival
priority weight receiving deadline credit. Completion-conditioned latency
and throughput can describe a different subset of requests. A policy that
finishes only favorable requests can have a strong completion-conditioned
metric while losing arrival-normalized utility; the common denominator
and missing-value rules are therefore part of the comparison, not merely
reporting details.

Across the complete two-pass campaign, 211 of 9,360 unique configurations
have zero completions; 48 of these belong to the 11-policy primary panel.
The preserved FIFO calibration manifest instead has zero zero-completion records out of 720. Calibration records cannot be propagated as zero-completion
outcomes for all other policies. The reported rankings retain defined
policy/window observations and omit undefined conditional observations,
so effective observation sets can differ even when a condition-level
ranking remains estimable.

The full per-policy outcome matrix remains the record for within-region
performance. This paper does not turn the panel into a universal winner
recommendation: the policies share a simulator implementation boundary,
and region names, synthesized deadlines and metric conditioning all
constrain the interpretation of an ordering.

\subsection{Absolute utility before rank comparison}
Table~\ref{tab:within-values} is a new descriptive tabulation of the frozen
outcomes, not a new execution or a replacement estimate. It reports the
mean ANWG of FIFO, the range of policy means, and the size of the exact
maximizing set. Each mean has the same 40-window support. Rounding in the
table does not determine ties; sets are computed from unrounded values.

\begin{table}[t]\centering
\caption{Within-source ANWG context from frozen outcomes. Minimum and maximum
are over the 11 policy means, not confidence limits. Top-set size uses exact
stored arithmetic. A=Azure, B=Bailian, G=BurstGPT.}
\label{tab:within-values}
\begin{tabular}{llrrrr}\toprule
Source & Region & FIFO & Minimum & Maximum & Top-set size\\\midrule
A & LOW & 1.000 & 0.018 & 1.000 & 6\\
A & PRE\_KNEE & 1.000 & 0.018 & 1.000 & 5\\
A & KNEE & 0.998 & 0.018 & 1.000 & 4\\
A & POST\_KNEE & 0.991 & 0.018 & 1.000 & 4\\
A & OVERLOAD & 0.987 & 0.018 & 1.000 & 4\\
A & HIGH\_PRESSURE & 0.964 & 0.018 & 1.000 & 1\\
B & LOW & 1.000 & 0.080 & 1.000 & 7\\
B & PRE\_KNEE & 1.000 & 0.080 & 1.000 & 6\\
B & KNEE & 0.998 & 0.080 & 1.000 & 4\\
B & POST\_KNEE & 0.988 & 0.080 & 1.000 & 4\\
B & OVERLOAD & 0.979 & 0.080 & 1.000 & 4\\
B & HIGH\_PRESSURE & 0.946 & 0.079 & 1.000 & 1\\
G & LOW & 1.000 & 0.968 & 1.000 & 9\\
G & PRE\_KNEE & 1.000 & 0.856 & 1.000 & 7\\
G & KNEE & 0.997 & 0.666 & 0.997 & 4\\
G & POST\_KNEE & 0.981 & 0.573 & 0.981 & 3\\
G & OVERLOAD & 0.935 & 0.490 & 0.936 & 2\\
G & HIGH\_PRESSURE & 0.645 & 0.271 & 0.652 & 2\\
\bottomrule\end{tabular}\end{table}
At LOW the maximum mean is 1.000 for every source, with six, seven and
nine policies tied in Azure, Bailian and BurstGPT respectively. Yet the
minimum is 0.018 for Azure, 0.080 for Bailian and 0.968 for BurstGPT.
A ranking difference at LOW therefore has quite different absolute-utility
context across sources. The table also explains why a tie-broken top
choice is a weak description of a unique preferred policy in this region.

At HIGH\_PRESSURE, FIFO means are 0.964, 0.946 and 0.645 respectively.
The highest means remain 1.000 for Azure and Bailian, whereas BurstGPT's
maximum is approximately 0.652. Thus similar ordinal agreement can coexist
with different attained utility. A high cross-source tau does not establish
similar service quality; a lower tau does not identify the magnitude or
cost of every changed pair. This is why Q1's ordinal summary is accompanied
by Q2's practical-margin criterion rather than interpreted alone.

These contrasts cannot be attributed solely to queue ordering. Classical
and serving-inspired implementations have different historical service
and resource semantics (Section~\ref{sec:execution-model}). In particular,
instant-prefill and prefill-aware paths do not impose the same work model.
Policies without hard KV enforcement can also report demand above nominal
capacity. The table supports description of the resulting abstractions;
it does not show that one native serving mechanism dominates another.
This is a fundamental comparability boundary for a scheduler-selection
claim, not an uncertainty that more resampling can remove.

\subsection{Observation support and operating-region interpretation}
All 48 zero-completion primary configurations belong to the historical
vLLM policy: seven Azure windows and one Bailian window in each of six
regions. Other primary policies provide defined completion outcomes on
those same inputs. Consequently, the corresponding conditional vLLM means
have 33 Azure or 39 Bailian observations, while other policy means use 40.
ANWG includes all 40 windows and gives the zero-completion outcomes zero
credit. These are different questions about the same frozen executions.

The zero-completion pattern is not a finding that each of those windows
is intrinsically infeasible for all policies. It is specific to an
implementation and its model path. Similarly, FIFO-referenced load factors
are controlled arrival-rate settings, not measured native saturation
states shared by every implementation. Interpreting higher-region changes
requires both the absolute outcomes and these implementation constraints.
The corpus does not estimate a universal utilization knee or a calibrated
capacity curve for production hardware.
